\graphicspath{{content/chapters/2_background/figures/}}

\chapter{Background (15 marks)}%
\label{chp:background}

% | Section                     | What Goes Here                             |
% | --------------------------- | ------------------------------------------ |
% | 2.1 Vehicle Routing Problem | Formal VRP definition                      |
% | 2.2 Dynamic VRP             | Dynamic arrivals / online routing          |
% | 2.3 Time Windows            | Constraints definition                     |
% | 2.4 DVRPTW Formulation      | Mathematical framing / objective           |
% | 2.5 Solomon Benchmark       | Dataset structure, benchmark conventions   |
% | 2.6 Evaluation Metrics      | Distance, vehicles used, runtime, lateness |



% How well and clearly did you:
% - Demonstrate reasoning and judgment for inclusion of relevant topics in the background;
% - Demonstrate clear narrative flow leading to the identification of the problem or focus area;
% - Define domain-specific concepts, methods and techniques used in the domain;
% - Manage to evidently and smoothly move from generic concepts and notions to include more
% specific detail and therefore focus



% Target Audience -> Academical and technical readers with some background in comupter science
% Research Question -> DCVRPTW: Heuristics vs Hybrid vs AI on Solomon C-RC-100 benchmark



% When writing the background for a study, start by providing a brief overview of the research topic 
% and its significance in the field.
% Then, highlight the gaps in existing knowledge or unresolved issues that the study aims to address. 
% Finally, summarize the key findings from relevant literature to establish the context and rationale 
% for conducting the research, emphasizing the need and importance of the study within the broader academic 
% landscape.

% Intro
Logistics optimisation concerns the systematic allocation and coordination of 
transportation, inventory, and distribution resources to minimise operational 
costs whilst maximising service efficiency within supply chain networks.
The importance of logistics optimisation has grown significantly in recent years, 
driven by the expansion of e-commerce platforms and increasing consumer expectations for rapid, on-demand delivery fulfilment.
Many logistical decision-making problems can be mathematically formulated as combinatorial optimisation tasks, 
wherein the objective is to determine the best solution from a finite but often exponentially large set of possible alternatives.
A prominent example is the Vehicle Routing Problem (VRP), an NP-hard combinatorial optimisation problem concerned with determining 
cost-efficient routes for a fleet of vehicles to service customers whilst satisfying practical constraints.

% VRP
\section{Vehicle Routing Problem (VRP)} % Adapted from RouteFinder paper, I don't know if I have to cite it here or not
The vehicle routing problem (VRP) can be defined on a graph $G$ as follows:
$G=(N,E)$, where $N = N_d \cup N_c$ is the set of nodes,
$N_d=\{0,\dots,m-1\}$ is the set of depots, and
$N_c=\{m,\dots,m+n-1\}$ is the set of customers.
Each edge $(i,j)\in E$ is associated with a travel cost $c_{ij}$.
A fleet of vehicles departs from the depot(s), serves each customer exactly once,
and minimizes the total travel cost.
Let $x_{ijk}\in\{0,1\}$ indicate whether vehicle $k$ traverses edge $(i,j)$.
The objective is to minimize the total travel cost:
\[
\min \sum_{k\in K}\sum_{i\in N}\sum_{j\in N\setminus\{i\}} c_{ij}x_{ijk}.
\]
The routing decisions satisfy
\begin{itemize}
	\item Assignment Constraint: Each customer is served by exactly one vehicle.
\[
\sum_{k\in K}\sum_{j\in N\setminus\{i\}} x_{ijk}=1,
\qquad \forall i\in N_c,
\]
	\item Flow Conservation: Each vehicle that arrives at a customer must also depart from it.
\[
\sum_{j\in N\setminus\{i\}} x_{ijk}
=
\sum_{j\in N\setminus\{i\}} x_{jik},
\qquad \forall i\in N_c,\; \forall k\in K.
\]
	\item Vehicle Returns to Depot: Each vehicle must start and end at the depot.
\[
\sum_{j\in N\setminus\{0\}} x_{0jk} = 1,
\qquad \forall k\in K,
\]
\[
\sum_{i\in N\setminus\{0\}} x_{i0k} = 1,
\qquad \forall k\in K.
\]
\end{itemize}


% Variants
\section{VRP Variants}
To model more realistic delivery environments, several extensions of the traditional VRP have been proposed to incorporate additional operational constraints.
\paragraph{CVRP}
The Capacitated Vehicle Routing Problem (CVRP) extends the classical VRP by introducing vehicle capacity constraints,
in the form of a maximum load that each vehicle can carry, $Q$,  and a demand $q\in[0,Q]$ associated with each customer 
such that the demand of each customer does not exceed the vehicle capacity, i.e., $q_i \leq Q$ for all $i\in N_c$.
Formally, Capacity Constraint can be expressed as follows:
The total demand assigned to a vehicle must not exceed its capacity:
\[
\sum_{i\in N_c} q_i y_{ik} \leq Q,
\qquad \forall k\in K where $K$ denotes the set of vehicles.
\]
\paragraph{VRPTW}
In addition to the constraints of CVRP, VRPTW introduces time window constraints for each customer $i\in N_c$,
defined by a time interval $[a_i, b_i]$ within which the service must start, often representing business operating
hours or delivery appointment slots.
Moreover, each customer now has a service time $s_i$ that represents the time required to complete the service at customer $i$.
Formally, the VRPTW additional constraints can be expressed as follows:
\begin{itemize}
	\item Time Window Constraint: Each customer must be served within their specified time window.
	\[
	a_i \leq t_i \leq b_i,
	\qquad \forall i\in N_c.
	\]
	\item Service Time Constraint: The service time at each customer must be accounted for in the scheduling of routes.
	Let $\tau_{ij}$ denote the travel time from node $i$ to node $j$.
	\[
t_j \geq t_i + s_i + \tau_{ij} - M(1 - x_{ijk}),
\qquad \forall i,j\in N_c,\; \forall k\in K,
\]
where $M$ is a sufficiently large constant to ensure the constraint is only active when vehicle $k$ travels from customer $i$ to customer $j$.
	
\end{itemize}

\paragraph{DVRPTW}
A more advanced extension is the Dynamic Capacitated Vehicle Routing Problem with Time Windows (DCVRPTW), 
which incorporates all aforementioned constraints whilst accounting for dynamic and uncertain events that may occur during operation. 
In DVRPTW, routing decisions must be continuously adapted as new customer requests arrive over time. 
This requires algorithms to operate in an online setting, where only partial information is available at any given moment, and previously planned routes may need to be modified dynamically.
Instead of a fixed customer set $N_c$ known in advance, we define: 
$N_c(t)$ as the set of customers known at time $t$, which can change dynamically. 
\paragraph{Importance of DVRPTW}
The DVRPTW is particularly relevant in real-world logistics systems, where customer requests may arrive dynamically throughout the day. 
Applications include same-day delivery, ride-sharing services, and emergency response systems, where decisions must be made in real time under uncertainty. 
The dynamic nature of the problem significantly increases its complexity, as routing plans must be continuously updated while maintaining service quality and operational efficiency.
\section{Benchmark Datasets}
Due to the importance of VRPTW, standardised benchmark datasets have been developed to ensure evaluations are fair and comparable. 
One of the most important for VRPTW is the Solomon benchmark, which consists of 100 instances with 100 customers each, designed to test the performance of VRPTW algorithms under different conditions. \cite{solomonAlgorithmsVehicleRouting1987}
\subsection{Geographical Distributions}
The dataset is split into three classes based on the spatial arrangement of customers:
\begin{itemize}
\item C (Clustered): Customers are grouped in distinct clusters, which typically results in shorter total distances
\item R (Random): Customers are uniformly distributed across the map; these are widely considered the most difficult instances because they lack clear structural patterns
\item RC (Mixed): A combination of clustered and random distributions
\end{itemize}

\subsection{Temporal and Capacity Constraints (Type 1 vs. Type 2)}
Solomon further divided these into two categories based on the tightness of constraints:
\begin{itemize}
	\item Type 1 (e.g., C1, R1, RC1): These have a short scheduling horizon, narrow time windows, and small vehicle capacities, which typically require multiple vehicles to service customers and result in more complex routing decisions.
	\item Type 2 (e.g., C2, R2, RC2): These have a long scheduling horizon, wide time windows, and large vehicle capacities, allowing many customers to be serviced by a single vehicle, which can simplify routing decisions but may also lead to longer routes and increased total distance.
\end{itemize}

\subsection{Adaptation for Dynamic VRP (DVRPTW)}
Although the original Solomon benchmark is static, it has been widely adapted for dynamic VRP research by introducing an "appearance time" for a percentage of the customers based on a Degree of Dynamism (DoD) $\epsilon$.
For example, in an instance with 50\% dynamism, half of the customers are hidden at the start of the day and only revealed as the vehicles are already executing their routes, simulating real-world scenarios 
where new orders can arrive unpredictably throughout the day.
\subsection{Common Preprocessing Steps}
% In your implementation section, note that it is standard practice to scale the raw Solomon coordinates—which usually occupy a 
% square—down to a 
% unit square
% This normalization is crucial for the stability of neural network training in models like the Attention Model (AM) or RouteFinder

\section{Evaluation Metrics}
The performance of algorithms addressing the Dynamic Vehicle Routing Problem with Time Windows (DVRPTW) 
is typically assessed through a range of evaluation metrics, each capturing different 
operational and computational aspects of solution quality. Since DVRPTW is inherently 
a multi-objective optimisation problem, no single metric is sufficient to fully 
describe algorithmic performance. Consequently, prior literature evaluates methods 
according to service quality, routing efficiency, computational efficiency, 
and decision-making effectiveness. The following metrics represent the most commonly 
adopted measures in recent DVRPTW and related dynamic routing studies.
\subsection{Service Quality Metrics}
\begin{itemize}
	
	\item \cite{chen2026} Lateness: Lateness is defined as the sum of delays beyond the latest allowable service time, typically expressed as $\sum_i \max(0, t_i - b_i)$.
	\item Rejection rate: Percentage of requests that are rejected due to capacity or time window constraints.% \cite{AntColony}
	\item Responsiveness: It measures how quickly a routing algorithm adapts to newly arriving requests, typically quantified as the delay between request arrival and its service start time. \cite{AntColony} 
\end{itemize}
\subsection{Operational Efficiency Metrics}
\begin{itemize}
	\item Total Distance: The Euclidean sum of all arcs traversed by the entire fleet. %Ojeda Rios et al. (2021) Section 3.4
	\item Fleet Size: The total number of vehicles (K or V) required to fulfill all demands. %Shahbazian et al. (2025) Table 4
	\item Total Duration: Cumulative time from depot departure until the last vehicle returns. %-Xu et al. (2025)
\end{itemize}

\subsection{Computational Efficiency Metrics}
\begin{itemize}
	\item Solving Time: The wall-clock computation time (s) required to return a solution for a snapshot. %Dornemann (2023) Table 3
\end{itemize}

\subsection{Benchmarking and Comparative Analysis}
\begin{itemize}
	\item Optimality Gap: The percentage difference between the algorithm's cost and the Best Known Solution. %Berto et al. (2024) Table 1
	\item Value of Information: The gap between real-time dynamic results and an "Oracle" (offline optimal) solution that has all the knowledge beforehand.\cite{valueOfInformation} 
\end{itemize}
\section{Solution Methods}
The VRP is NP-hard, making exact methods computationally infeasible for large instances. To address this, researchers have
developed various heuristic and metaheuristic approaches that provide good quality solutions within reasonable time frames.
\subsection{Heuristics}
 Heuristics are problem-specific strategies that construct solutions based on intuitive rules or domain knowledge. These can be broadly categorised into two types:
\begin{itemize}
	\item{Constructive heuristics}: These build a solution from scratch by iteratively adding customers to routes based on criteria such as proximity or cost. Examples include the Nearest Neighbor, Clarke-Wright Savings, and Solomon's Insertion heuristics.\cite{metaheuristicsAndHeuristics}
	\item {Improvement heuristics}: These start with an initial solution and iteratively improve it through local search techniques, such as 2-opt, 3-opt, or Or-opt, which involve swapping or relocating customers to reduce total distance or improve route feasibility.\cite{metaheuristicsAndHeuristics}
\end{itemize}
\subsection{Metaheuristics}
 Metaheuristics are higher-level frameworks that guide the search process to explore the solution space more effectively.
They often incorporate mechanisms to escape local optima and can be applied to a wide range of combinatorial problems. \cite{metaheuristicsAndHeuristics}
Common metaheuristics for VRP include:
\begin{itemize}
	\item Genetic Algorithms (GA): Use evolutionary principles to evolve a population of solutions through selection, crossover, and mutation.
	\item Tabu Search(TS): Employs memory structures to avoid cycling back to previously visited solutions.
	\item Guided Local Search(GLS): Enhances local search by penalizing certain solution features to escape local minima.
	\item Simulated Annealing (SA): Mimics the cooling process of metals to probabilistically accept worse solutions early on to escape local minima.
	\item Ant Colony Optimization (ACO): Inspired by the foraging behavior of ants, it uses pheromone trails to guide the search for optimal routes.
\end{itemize}

\section{AI Methods}
In addition to classical heuristics, recent work has explored AI-based approaches for solving combinatorial optimisation problems such as the VRP. 
These methods aim to learn solution construction strategies directly from data, rather than relying on manually designed rules. 
To enable this, optimisation problems are reformulated in a way that allows them to be processed by neural network models.

A fundamental requirement in this setting is the representation of problem instances. 
Routing problems are naturally expressed as graphs, where nodes correspond to locations and edges encode travel costs or distances. 
Each node is associated with a feature representation that captures relevant problem attributes and enables processing by learning-based models.

To capture interactions between nodes, attention mechanisms are commonly employed. 
Attention mechanisms allow models to dynamically focus on relevant parts of the input when constructing solutions, enabling both local and global relationships between nodes to be captured.

Solution construction is typically formulated as a sequential process. 
A solution is generated step-by-step, where at each step the model selects the next node based on the current partial solution and the encoded problem representation. 
This is achieved through a decoder that produces a probability distribution over feasible actions, while enforcing constraints such as visiting each customer only once. 
As a result, the complete solution can be viewed as a sequence of decisions.

To learn such decision-making strategies, reinforcement learning is commonly used. 
In this framework, the quality of a constructed solution is evaluated using a reward signal, typically defined in terms of routing cost, and the model is trained to maximise expected reward over a distribution of problem instances. 
This allows the model to learn effective solution strategies without requiring optimal solutions as supervision \cite{Kool2019}.

% These concepts, graph-based representations, attention-based encoding, and sequential decision-making, form the foundation of modern neural approaches to routing problems, and provide the necessary background to understand attention-based models such as the Attention Model and its extensions.